\documentclass[letterpaper,11pt]{article}

\usepackage[empty]{fullpage}
\usepackage[hidelinks]{hyperref}
\usepackage[usenames,dvipsnames]{color}
\usepackage[english]{babel}
\usepackage[right=2.2cm, top=2cm, left=2.2cm, bottom=2cm]{geometry}
\usepackage{titlesec}
\input{glyphtounicode}
\pdfgentounicode=1

\urlstyle{same}
\raggedbottom
\setlength{\parindent}{0pt}
\setlength{\parskip}{10pt}

\begin{document}

%----------HEADING (matches resume)----------
\begin{center}
  \textbf{\Huge \scshape Pablo Leyva} \\ \vspace{3pt}
  \small 908-525-6062 $|$ \href{mailto:pl33@njit.edu}{\underline{pl33@njit.edu}} $|$
  \href{https://linkedin.com/in/pablo-leyva}{\underline{linkedin.com/in/pablo-leyva}} $|$
  \href{https://github.com/pleyva2004}{\underline{github.com/pleyva2004}} $|$
  \href{https://pabloleyva.io}{\underline{pabloleyva.io}}
\end{center}

\vspace{6pt}
\hrule
\vspace{14pt}

Dear Google AI Career Catalyst Recruiting Team,

I am a senior at the New Jersey Institute of Technology completing a B.S. in Applied Mathematics and Computer Science (AI concentration, May 2027), and I am applying to the AI Career Catalyst Program because its premise, engineering at the intersection of AI and infrastructure, is exactly the intersection where I have spent the last three years building.

On the AI side, I am currently in my second internship at Apple, where I build model-based tools that curate and score RL post-training data for a continually learning small language model, and where I am implementing GRPO/PPO on-policy training with vLLM rollouts against verifiable-reward environments. Last summer at Apple I built an agent orchestration system on the Apple Foundation Model with MCP and the team's LLM evaluation pipeline. At NJIT, my research fellowship has me training a 3B-parameter reasoning model with multi-GPU distributed training (FSDP) through an SFT, DPO, and RL fine-tuning pipeline.

On the infrastructure side, I built NeuroCache, an open-source KV-cache memory manager with attention-scored eviction and cache-replacement policies, and at Caterpillar I built an Airflow pipeline ingesting 500K+ daily records with FastAPI model serving at sub-100ms latency. The models I care about only matter when the systems underneath them are efficient and reliable, and that is the layer I want to master at Google's scale.

The rotational format is why I am drawn to this program specifically: I have proven I ramp quickly in unfamiliar territory. I have interned across heavy industry and consumer technology, and with a team of five engineers at SHPE NJIT I built and shipped our chapter's first iOS app in 27 days, growing it to 200+ users in its first month. As SHPE chapter president I lead a 20-person board serving 300+ students. I grew up in Mexico City in a family that treated education as the way forward, and I bring that same conviction, and the responsibility to open doors for others, to every team I join.

I graduate in May 2027 and am available to start in June 2027. I would welcome the chance to speak further about the program. Thank you for your consideration.

Sincerely,

Pablo Leyva

\end{document}
